\subsection*{What a Hermitian generator fixes}
Work on a finite graph and finite-dimensional Hilbert space with
ruler operators $\delta I\le s_x\le MI$, $\delta>0$,
central real lapses, and a common
differentiable operator domain.
Requiring H2's row to be $C_x=\partial_{N_x}H$ for one Hermitian
generator fixes Hermiticity and the compatibility of mixed
derivatives.  It does not determine a unique factor ordering.

One explicit counterfamily starts from the scalar polar densities
$t_x(\psi)$, whose sum is zero, and chooses a linear
adjoint-preserving lift $T_x=\Phi[t_x;s]$.
Thus $\sum_xT_x=0$ is an operator identity.
For $\alpha\in\mathbb R$ set
\[
 \begin{aligned}
 \mathcal M_{s,\alpha}(Y)&=\tfrac12\{s,Y\}+i\alpha[s,Y],\\
 H[L;s]&=\sum_x\left[
 \mathcal M_{L_x,\alpha}(T_x)+\sum_rQ_{rx}^\dagger L_xQ_{rx}
 \right].
 \end{aligned}
\]
All adjoints are Hilbert-space adjoints.
The matter columns, including their excitation filter, are chosen
independently of the varied lapse.
For $L_x=N_xs_x$ the derived row is
$C_x=\mathcal M_{s_x,\alpha}(T_x)+\mathcal E_x$, where
$\mathcal E_x=\sum_rQ_{rx}^\dagger s_xQ_{rx}\ge0$.
The allocation is Hermitian and affine for every $\alpha$, and its
vacuum part vanishes when all $L_x$ are the same operator.

The inverse of the ruler multiplier is positive for every $\alpha$:
\[
 \begin{aligned}
 \mathcal R_{s,\alpha}(Y)
 &=\int_0^\infty e^{-u(1/2+i\alpha)s}
       Y e^{-u(1/2-i\alpha)s}\,du,\\
 \mathcal M_{s,\alpha}\mathcal R_{s,\alpha}&=1.
 \end{aligned}
\]
Differentiating the integrand proves the inverse identity; each
integrand is a congruence, so $\mathcal R_{s,\alpha}$ is completely
positive and has norm at most $\delta^{-1}$.
In an eigenbasis of $s$ an independent check gives
\[
 (\mathcal R_{s,\alpha}Y)_{ij}
 =\frac{Y_{ij}}{(\lambda_i+\lambda_j)/2+
                    i\alpha(\lambda_i-\lambda_j)} .
\]
[X: convergent Lyapunov integral; spectral matrix-element solution.]
On strong operator constraints, or a physical subspace reducing all
rulers and the restricted construction,
\[
 -T_x=\mathcal R_{s_x,\alpha}(\mathcal E_x)\ge0.
\]
At a unit collar the same-law charge is
$B=-\sum_I T_x+\sum_{\rm collar}\mathcal E_x\ge0$.
The weaker condition $C_xP=0$ alone does not justify applying this
inverse on $P$.

Hermiticity and agreement with $LT$ for commuting factors classify
every constant-coefficient bilinear $aLT+bTL$:
$b=\overline a$ and $a+b=1$, giving exactly this $\alpha$ family.
Its distinct denominators can change the charge spectrum for
$[s,\mathcal E]\ne0$, as the explicit rank-1 example below verifies.
Requiring antiunitary time reversal with even $s,T$ selects
$\alpha=0$ within this class, since it sends $\alpha$ to $-\alpha$.
That is an additional selecting principle; multiruler lifts and
matter-column orderings remain unspecified.
A Hermitian row also does not imply a positive source:
for $s=\operatorname{diag}(1,2)$ and
$Y=\tfrac12(1,1)^T(1,1)$, $\{s,Y\}/2$ has eigenvalues
$(3\pm\sqrt{10})/4$.
The actual Gram source avoids that failure.
For the same $s,Y$, the inverse charge has trace $3/4$ and determinant
$1/8-[4(9/4+\alpha^2)]^{-1}$, so its spectrum varies with $|\alpha|$.

\paragraph*{SOURCES.}
\begingroup\footnotesize
\path{<internal>/P26_operator_theory.md},
actual attachment read completely, especially sections 1--6 and 9:
trace-first lift, the Hermitian allocation family, positive inverse,
strong-versus-weak constraint domain and complete source derivative.
For example $\Phi$ may be the average of a fixed ordered
multivariable resolvent calculus and its reverse on a compact
positive spectral box.  Linearity preserves the traced flux identity;
positivity of arbitrary scalar functions under this lift is not assumed.
\path{<internal>/STATE.tex},
the source/bridge [S] operator entry and open problem O5:
the constructed finite-domain family does not establish a physical
stellar representation, a scalar-limit error bound or its time law.
No numerical checks reported by P26 are treated as executions here.
\endgroup
% Governing source: <internal>/GOAL_alt.md,
% duty 4(4). Scalar traces must be formed before this chosen lift.
% A full trace of |A^dagger|-|A| is zero, but a partial trace is not
% cyclic for operator-valued entries. An explicit generic witness:
% on C^2 tensor C^2 take E=I tensor diag(1,2), U=SWAP, A=UE.
% Then E0=E, E1=diag(1,2) tensor I, and the partial trace over the
% first factor of E1-E0 is diag(1,-1), although the full trace is zero.
% This witnesses the invalid inference, not a special H_* polar stencil.
% For s=diag(1,2), E=(1,1)^T(1,1)/2, the inverse charge has trace
% 3/4 and determinant 1/8-1/[4(9/4+alpha^2)], so its spectrum
% depends on alpha even while its positivity is retained.
% At a collar s_b=I and N_b=1, the vacuum term is T_b and the
% matter contribution is E_b; retaining it yields the displayed B.
% Uniqueness of a strong representation and the open-exterior limit
% are not implied by the positive inverse formula.
